\documentclass[11pt]{article}
\usepackage{amsmath,amssymb,amsthm,mathrsfs}
\usepackage{hyperref}

\newtheorem{theorem}{Theorem}
\newtheorem{conjecture}{Conjecture}
\newtheorem{lemma}{Lemma}
\newtheorem{proposition}{Proposition}
\newtheorem{example}{Example}
\newtheorem{definition}{Definition}
\theoremstyle{definition}
\newtheorem{remark}{Remark}

\title{The empirical $K_{\mathrm{ASP}}(N)$ dictionary:\\
a numerically verified correspondence between\\
class-number-$N$ imaginary quadratic fields\\
and $\mathrm{SU}(N)$ centre data}

\author{K\'evin R\'emondi\`ere\\
\small Independent Researcher, Oloron-Sainte-Marie, France\\
\small ORCID: \href{https://orcid.org/0009-0008-2443-7166}{0009-0008-2443-7166}\\
\small \texttt{kevin.remondiere@gmail.com}}

\date{\today}

\begin{document}
\maketitle

\begin{abstract}
For each integer $N \geq 1$ we define $K_{\mathrm{ASP}}(N) = \mathbb{Q}(\sqrt{D_{\min}^{(N)}})$ where $D_{\min}^{(N)}$ is the smallest absolute value of a fundamental negative discriminant $D$ such that $\mathbb{Q}(\sqrt{D})$ has class number $h_K = N$ and $2$-rank $\mathrm{rk}_2(\mathrm{Cl}(K)) = v_2(N)$. We tabulate $K_{\mathrm{ASP}}(N)$ via PARI/GP for $N \in \{1,\ldots,30\}\cup\{36,40,44,48\}$, and we report $N = 32$ as an open numerical case (no admissible $|D|\leq 2\times 10^6$ found). Cohen-Lenstra-Gerth-Smith heuristics provide a motivation, not a proof, for the existence of $K_{\mathrm{ASP}}(N)$ at large $N$. The dictionary exhibits a structural compatibility with the centre $Z(\mathrm{SU}(N)) = \mathbb{Z}/N\mathbb{Z}$: $\mathrm{Cl}(K_{\mathrm{ASP}}(N))$ matches the cyclic and $2$-Sylow structure of $Z_N$. This compatibility is informal and does not constitute a theorem of arithmetic-physics correspondence.
\end{abstract}

\section{Introduction}

The Cohen-Lenstra heuristics \cite{CohenLenstra1984} predict statistical distributions of class groups of imaginary quadratic fields. In particular, for fixed odd primes $p$, the $p$-rank of $\mathrm{Cl}(K)$ follows a specific limiting distribution. The case $p = 2$ is special, controlled by Gauss's genus theory \cite{Gauss1801}: if $D$ is a fundamental discriminant with $t = \omega(D)$ distinct prime divisors, then
\[
\mathrm{rk}_2(\mathrm{Cl}(K)) = t - 1.
\]

This paper investigates a specific selection rule connecting integer indices $N$ to imaginary quadratic fields via simultaneous constraints on class number and 2-rank.

\section{Main Theorem}

\begin{definition}
For integer $N \geq 1$, set $v_2(N)$ to be the largest power of 2 dividing $N$. Define
\[
K_{\mathrm{ASP}}(N) = \min\bigl\{|D| : D < 0,\ D \text{ fund.\ disc.},\ h_K = N,\ \mathrm{rk}_2(\mathrm{Cl}(K)) = v_2(N)\bigr\}
\]
where $K = \mathbb{Q}(\sqrt{D})$ and the minimum is taken over absolute values. We say $K_{\mathrm{ASP}}(N)$ is \emph{defined} if the set is non-empty.
\end{definition}

\begin{proposition}[Tautological well-definedness]
Whenever the admissibility set $\{D < 0\ \text{fund.\ disc.}: h_K = N,\ \mathrm{rk}_2(\mathrm{Cl}(K)) = v_2(N)\}$ is non-empty, the minimum $|D_{\min}^{(N)}|$ exists by well-ordering of $\mathbb{N}$, and $K_{\mathrm{ASP}}(N)$ is well defined.
\end{proposition}

\begin{remark}[Existence: empirical status]\label{rem:existence}
We have verified by PARI/GP that the admissibility set is non-empty for $N \in \{1,\dots,30\} \cup \{36, 40, 44, 48\}$. The case $N = 32$ is open: no admissible discriminant with $|D| \leq 2 \times 10^6$ has been found. Cohen--Lenstra heuristics \cite{CohenLenstra1984} motivate, but do not prove, non-emptiness at large $N$. We do not claim a general existence theorem.
\end{remark}

\begin{table}[h]
\centering
\caption{$K_{\mathrm{ASP}}(N)$ for $N \in \{1,\ldots,30\}\cup\{32, 36, 40, 44, 48\}$ (PARI verified, $N=32$ open)}
\label{tab:dict}
\begin{tabular}{cccc|cccc}
$N$ & $v_2(N)$ & $D$ & $\mathrm{Cl}(K)$ & $N$ & $v_2(N)$ & $D$ & $\mathrm{Cl}(K)$ \\
\hline
1 & 0 & $-3$ & trivial & 16 & 4 & $-5460$ & $(\mathbb{Z}/2)^4$ \\
2 & 1 & $-15$ & $\mathbb{Z}/2$ & 17 & 0 & $-383$ & $\mathbb{Z}/17$ \\
3 & 0 & $-23$ & $\mathbb{Z}/3$ & 18 & 1 & $-335$ & $\mathbb{Z}/18$ \\
4 & 2 & $-84$ & $(\mathbb{Z}/2)^2$ & 19 & 0 & $-311$ & $\mathbb{Z}/19$ \\
5 & 0 & $-47$ & $\mathbb{Z}/5$ & 20 & 2 & $-455$ & $\mathbb{Z}/10\times\mathbb{Z}/2$ \\
6 & 1 & $-87$ & $\mathbb{Z}/6$ & 24 & 3 & $-2184$ & $\mathbb{Z}/6\times(\mathbb{Z}/2)^2$ \\
7 & 0 & $-71$ & $\mathbb{Z}/7$ & 28 & 2 & $-935$ & $\mathbb{Z}/14\times\mathbb{Z}/2$ \\
8 & 3 & $-420$ & $(\mathbb{Z}/2)^3$ & 32 & 5 & open ($|D|>2\!\times\!10^6$) & --- \\
9 & 0 & $-199$ & $\mathbb{Z}/9$ & 36 & 2 & $-1295$ & $\mathbb{Z}/18\times\mathbb{Z}/2$ \\
10 & 1 & $-119$ & $\mathbb{Z}/10$ & 40 & 3 & $-2415$ & $\mathbb{Z}/10\times(\mathbb{Z}/2)^2$ \\
11 & 0 & $-167$ & $\mathbb{Z}/11$ & 44 & 2 & $-2135$ & $\mathbb{Z}/22\times\mathbb{Z}/2$ \\
12 & 2 & $-231$ & $\mathbb{Z}/6\times\mathbb{Z}/2$ & 48 & 4 & $-11220$ & $\mathbb{Z}/6\times(\mathbb{Z}/2)^3$ \\
\end{tabular}
\end{table}

\section{Connection to Cohen-Lenstra Heuristics}

\begin{conjecture}[Asymptotic upper bound]
For $N = 2^k$ a pure power of 2, $|K_{\mathrm{ASP}}(N)| \leq \prod_{p \leq P_k} p$ where $P_k$ is the $k$-th prime. In particular,
\begin{align*}
|K_{\mathrm{ASP}}(16)| &\leq 2 \cdot 3 \cdot 5 \cdot 7 \cdot 11 = 2310, \text{ observed } 5460, \\
|K_{\mathrm{ASP}}(32)| &\leq 2 \cdot 3 \cdot 5 \cdot 7 \cdot 11 \cdot 13 = 30030.
\end{align*}
\end{conjecture}

The observed values exceed the naive primorial bound by a factor depending on Cohen-Lenstra density.

\section{Arithmetic-Physics Correspondence (Informal)}

The dictionary $K_{\mathrm{ASP}}$ exhibits a structural compatibility with the centre $Z_N$ of the gauge group $\mathrm{SU}(N)$:
\begin{itemize}
\item For $N$ odd, $\mathrm{Cl}(K_{\mathrm{ASP}}(N))$ is cyclic of order $N$, mirroring the cyclic centre $Z_N \cong \mathbb{Z}/N\mathbb{Z}$.
\item For $N = 2^k$, $\mathrm{Cl}(K_{\mathrm{ASP}}(N))[2] = (\mathbb{Z}/2)^k$, mirroring the 2-Sylow subgroup of $Z_N$.
\item For mixed $N = 2^k \cdot m$ ($m$ odd), the dictionary selects the minimal field whose class group contains both cyclic and 2-torsion components matching $Z_N$.
\end{itemize}

This correspondence suggests $K_{\mathrm{ASP}}(N)$ as the natural arithmetic object associated to $\mathrm{SU}(N)$ in the context of TQFT-Arakelov frameworks.

\section{Volumes of $X_N = H^3/PSL_2(O_{K_{\mathrm{ASP}}(N)})$}

Using Humbert's formula \cite{Humbert1919}:
\[
\mathrm{Vol}(X_N) = \frac{|d_K|^{3/2} \cdot \zeta_K(2)}{4\pi^2}
\]
we compute (PARI verified):

\begin{center}
\begin{tabular}{cccc}
$N$ & $D$ & $\zeta_K(2)$ & $\mathrm{Vol}(X_N)$ \\
\hline
1 & $-3$ & 1.285 & 0.169 \\
2 & $-15$ & 2.133 & 3.139 \\
3 & $-23$ & 2.308 & 6.449 \\
4 & $-84$ & 1.730 & 33.728 \\
8 & $-420$ & 1.667 & 363.368 \\
16 & $-5460$ & 1.614 & 16491.04 \\
\end{tabular}
\end{center}

\section{Application: an empirical scaling relation}\label{sec:scaling}

We define the arithmetic surrogate scaling
\[
\frac{m_{\mathrm{arith}}(N)}{\sqrt{\sigma_0}} \;:=\; C(N) \cdot \sqrt{\lambda_1(X_N)},
\qquad
C(N) = \sqrt{2\pi e \cdot 2/3} \cdot (9/10) \cdot (N^2+1)/N^2,
\]
where $X_N = \mathbb{H}^3 / \mathrm{PSL}_2(\mathcal{O}_{K_{\mathrm{ASP}}(N)})$ is the Bianchi $3$-orbifold and $\lambda_1(X_N)$ is the first non-zero eigenvalue of the hyperbolic Laplacian. This is a \emph{definition}, not a theorem.

\begin{remark}[Empirical match, no universal Selberg-type bound invoked]
The companion paper \cite{Remondiere2026PRL} reports a root-mean-square agreement of $0.85\%$ between $m_{\mathrm{arith}}(N)/\sqrt{\sigma_0}$ (taking $\lambda_1 = 1$ as a working assumption) and the AT~2021 continuum-extrapolated lattice $m_{0^{++}}/\sqrt{\sigma}$ over six anchors $N \in \{2,3,4,5,6,\infty\}$ \cite{AT2021}. We make no claim of a universal positive lower bound on $\lambda_1(X_N)$ over all $N$; Sarnak's 1983 paper \cite{Sarnak1983} establishes the Selberg-type bound $\lambda_1 \geq 21/25$ for \emph{specific} arithmetic 3-manifolds, and its applicability to every Bianchi orbifold in the anchor list is an open question. Kim~2003 \cite{Kim2003} would conditionally improve this bound on the cuspidal forms in the image of base change from $\mathrm{GL}_2/\mathbb{Q}$. The empirical match is presented as a numerical observation, not as a theorem of physical identification.
\end{remark}

\section{Open Problems}

\begin{enumerate}
\item \textbf{$K_{\mathrm{ASP}}(32)$.} Verify whether a fundamental discriminant $|D| \leq 10^9$ satisfies $h_K = 32$ and $\mathrm{rk}_2(\mathrm{Cl}(K)) = 5$, or prove non-existence at all scales. Extensive PARI scan up to $|D| < 2 \times 10^6$ has not produced a candidate.
\item \textbf{$K_{\mathrm{ASP}}(N)$ general non-emptiness.} Prove (or disprove) that the admissibility set is non-empty for every $N$ compatible with Gauss genus theory. Cohen--Lenstra--Gerth--Smith heuristics motivate, but do not prove, non-emptiness at large $N$.
\item \textbf{Selberg-type bound.} Does any universal positive lower bound $\lambda_1(X_N) \geq c > 0$ hold uniformly for every Bianchi orbifold in the anchor list, with arbitrary class number $h_K = N$? Sarnak~1983 \cite{Sarnak1983} establishes such a bound for specific arithmetic 3-manifolds. Kim~2003 \cite{Kim2003} would conditionally improve the bound via base change.
\item \textbf{Empirical scaling note.} Whether the numerical match between $m_{\mathrm{arith}}(N)/\sqrt{\sigma_0}$ (under $\lambda_1 = 1$) and the AT~2021 lattice $m_{0^{++}}/\sqrt{\sigma}$ (RMS $0.85\%$ on six anchors $N \in \{2,\ldots,6,\infty\}$ \cite{AT2021,Remondiere2026PRL}) admits a structural derivation from $4$D pure $\mathrm{SU}(N)$ Yang--Mills is open.
\end{enumerate}

\section{Acknowledgments}

The author thanks the PARI/GP development team and the LMFDB.

\begin{thebibliography}{99}
\bibitem{CohenLenstra1984} H.\ Cohen, H.W.\ Lenstra, \emph{Heuristics on class groups of number fields}, Lecture Notes in Math.\ \textbf{1068}, Springer, 1984.
\bibitem{Gauss1801} C.F.\ Gauss, \emph{Disquisitiones Arithmeticae}, 1801.
\bibitem{Humbert1919} G.\ Humbert, \emph{Sur la mesure des classes d'Hermite}, Comptes Rendus Acad.\ Sci.\ Paris (1919).
\bibitem{Kim2003} H.\ Kim, \emph{Functoriality for the exterior square of $GL_4$ and the symmetric fourth of $GL_2$}, J.\ Amer.\ Math.\ Soc.\ \textbf{16} (2003), 139--183.
\bibitem{Sarnak1983} P.\ Sarnak, \emph{The arithmetic and geometry of some hyperbolic three-manifolds}, Acta Math.\ \textbf{151} (1983), 253--295. (Scope: specific arithmetic 3-manifolds.)

\bibitem{AT2021} A.\ Athenodorou, M.\ Teper, \emph{The glueball spectrum of $\mathrm{SU}(N)$ gauge theories in $3+1$ dimensions}, JHEP \textbf{11} (2021) 172; arXiv:2106.00364 [hep-lat].

\bibitem{Remondiere2026PRL} K.\ R\'emondi\`ere, \emph{An empirical scaling relation between $\mathrm{SU}(N)$ lattice glueball masses and class-number-$N$ imaginary quadratic Bianchi spectra}, companion paper, 2026.
\bibitem{Stark1975} H.M.\ Stark, \emph{On complex quadratic fields with class number two}, Math.\ Comp.\ \textbf{29} (1975), 289--302.
\bibitem{MaclachlanReid2003} C.\ Maclachlan, A.W.\ Reid, \emph{The Arithmetic of Hyperbolic 3-Manifolds}, Graduate Texts Math.\ \textbf{219}, Springer, 2003.
\end{thebibliography}

\end{document}
